%% settings for listings-package %%
% basline listings style
\lstdefinestyle{BaselineStyle}
{ %
    backgroundcolor=\color{BlanchedAlmond!20},   
    commentstyle=\color{Blue},
    keywordstyle=\color{Green},
    numberstyle=\tiny\color{Gray},
    stringstyle=\color{Fuchsia},
    basicstyle=\ttfamily\scriptsize,
    columns=flexible,
    breakatwhitespace=false,         
    breaklines=true,                 
    captionpos=b,                    
    keepspaces=true,                 
    numbers=left,                    
    numbersep=5pt,                  
    showspaces=false,                
    showstringspaces=false,
    showtabs=false,                  
    tabsize=2,
    xleftmargin=2em,
    frame=single,
    framexleftmargin=1.5em
} %
% set to 'BaselineStyle'
\lstset{style=BaselineStyle}

% TigerJython syntax
\lstdefinestyle{TigerJython}
{ %
    language=Python,
    morekeywords={repeat},
    backgroundcolor=\color{Pink!15},   
} %

%% tikz %%
\usetikzlibrary{
  3d,
  arrows,
  arrows.spaced,
  arrows.meta,
  bending,
  babel,
  calc,
  fit,
  patterns,
  patterns.meta,
  plotmarks,
  shapes.geometric,
  shapes.misc,
  shapes.symbols,
  shapes.arrows,
  shapes.callouts,
  shapes.multipart,
  shapes.gates.logic.US,
  shapes.gates.logic.IEC,
  circuits.logic.US,
  circuits.logic.IEC,
  circuits.logic.CDH,
  circuits.ee.IEC,
  datavisualization,
  datavisualization.polar,
  datavisualization.formats.functions,
  er,
  automata,
  backgrounds,
  chains,
  topaths,
  trees,
  petri,
  mindmap,
  matrix,
  calendar,
  folding,
  fadings,
  shadings,
  spy,
  through,
  turtle,
  positioning,
  scopes,
  decorations.fractals,
  decorations.shapes,
  decorations.text,
  decorations.pathmorphing,
  decorations.pathreplacing,
  decorations.footprints,
  decorations.markings,
  shadows,
  lindenmayersystems,
  intersections,
  fixedpointarithmetic,
  fpu,
  svg.path,
  external,
  graphs,
  graphs.standard,
  quotes,
  math,
  angles,
  views,
  animations,
  rdf,
  perspective,
  tikzmark,
  optics,
  overlay-beamer-styles
}

\tikzstyle{code} = [rectangle, 
rounded corners, 
minimum width=3cm, 
minimum height=1cm,
text centered, 
draw=black, 
fill=red!30]

\tikzstyle{decision} = [diamond,
aspect=2,
minimum width=3cm, 
minimum height=1cm, 
text centered, 
draw=black,
fill=green!30]
\tikzstyle{arrow} = [thick,->,>=stealth]
\tikzset{>=latex}
\tikzset{boximg/.style={remember picture,red,thick,draw,inner sep=0pt,outer sep=0pt}}
\newcommand\Ceasar[1][3]{%
	\begin{tikzpicture}
		\def\R{2}
		\def\dR{0.5}
		\def\Ri{1.4}
		\def\dRi{0.5}
		\def\Rotate{#1}%% Number of shifts
		\coordinate (origin) at (0,0);
		\draw[fill=gray!60] (origin) circle (1pt);
		\draw (origin) circle (\R);
		\draw (origin) circle (\R+\dR);
		\draw (origin) circle (\Ri);
		\draw (origin) circle (\Ri+\dRi);
		\foreach \Letter [count=\ind from 0,evaluate=\ind as \ang using 90-\ind*360/26] in {A,B,...,Z}{%
			\node[rotate={\ang-90}] at ($(origin)+(\ang:{\R+\dR/2})$) {\Letter};
			\draw({\ang+0.5*360/26}:\R)--({\ang+0.5*360/26}:\R+\dR);
		}
		\foreach \Letter [count=\ind from 0,evaluate=\ind as \ang using 90-\ind*360/26-\Rotate*360/26] in {A,B,...,Z}{%
			\node[rotate={\ang-90}] at ($(origin)+(\ang:{\Ri+\dRi/2})$) {\Letter};
			\draw({\ang+0.5*360/26}:\Ri)--({\ang+0.5*360/26}:\Ri+\dRi);
		}
		\draw[-latex] (90:\Ri-0.3) arc (90:{90-\Rotate*360/26}:\Ri-0.3)node[pos=0.5,anchor=90-\Rotate*180/26]{$\Rotate$};
	\end{tikzpicture}
}


%% specific settings for different document classes %%
\iftoggle{exam}{
    %% EXAM ENVIRONMENT %%
    \pointpoints{Punkt}{Punkte}
    \bonuspointpoints{Bonuspunkt}{Bonuspunkte}
    \renewcommand{\solutiontitle}{\noindent\textbf{Lösungsvorschlag:}\enspace}
    \hqword{Aufgabe:}
    \hpgword{Seite:}
    \hpword{Punkte:}
    \hsword{erhalten:}
    \htword{Total} 
    \bhqword{Aufgabe:}
    \bhpgword{Seite:}
    \bhpword{Bonuspunkte:}
    \bhsword{erhalten:}
    \bhtword{Total} 
    \chqword{Aufgabe:}
    \chpgword{Seite:}
    \chpword{Punkte:}
    \chbpword{Bonuspunkte:}
    \chsword{erhalten:}
    \chtword{Total} 
    %\checkedchar{\CheckedBox}
    \totalformat{Aufgabe \thequestion\ total: \totalpoints\ Punkte}
    \runningheadrule
    %\firstpageheader{\mycourse}{\myclass}{\mydate}
    \runningheader{\mycourse}{\myclass}{\mydate}
    \firstpagefooter{}{\thepage\,/\,\numpages}{}
    \runningfooter{}{\thepage\,/\,\numpages}{}
    \CorrectChoiceEmphasis{\normalfont}
    % \renewcommand{\questionlabel}{\bfseries Aufgabe~\thequestion.}
}{}

\iftoggle{book}{
    % BOOK ENVIRONMENT %
    \setcounter{secnumdepth}{5}
    \setcounter{tocdepth}{5}
    % start chapter numbering at 0
    \setcounter{chapter}{-1}
    % slightly changing chapter marking in header
    \renewcommand{\chaptermark}[1]{\markboth{\textnormal{\thechapter}\ \textnormal{#1}}{}}
    % slightly changing section marking in header
    \renewcommand{\sectionmark}[1]{\markright{\textnormal{\thesection}\ \textnormal{#1}}{}}
    % making header of toc non italics (\upshape) and not all capitals
    \addto\captionsgerman{\renewcommand{\contentsname}{\upshape{I\MakeLowercase{nhaltsverzeichnis}}}}
}{}


% BEAMER %
\iftoggle{beamer}{
\beamertemplatenavigationsymbolsempty
}{}

%% replace English terms with German counterparts %%
\renewcommand{\lstlistingname}{Programm}
\floatname{algorithm}{Algorithmus}
\renewcommand{\algorithmicrequire}{\textbf{Eingabe:}}
\renewcommand{\algorithmicensure}{\textbf{Resultat:}}

%% clever refs %%
\nottoggle{beamer} {
    \crefname{lstlisting}{Programm}{Programme}
    \Crefname{lstlisting}{Programm}{Programme}
}{}